\chapter{Asynchronous Python: The Execution Model}
\label{ch:asyncio}

\begin{quote}\itshape
A coroutine is cold. That one sentence accounts for half the bugs.
\end{quote}

\noindent
If you write \code{async}/\code{await} in C\#, you already hold the concepts.
What you do not hold are Python's defaults, and they differ in ways that turn
correct C\# instincts into wrong Python code. A \code{Task} in .NET is usually
already running. A coroutine in Python has not started, and will not start until
something awaits it or hands it to the loop. There is one event loop per thread,
no synchronisation context, and no \code{ConfigureAwait} to reason about. And one
blocking call anywhere stops the entire process, not one request.

This chapter earns its claims with numbers rather than assertions. Every
measurement in it was run on the machine that produced this book, costs nothing,
needs no provider, and takes under a minute to reproduce.

% ============================================================
\section{The event loop}
\label{sec:event-loop}
\index{asyncio!event loop}
\index{concurrency!cooperative}

One thread. One loop. It runs one piece of your code at a time and switches
between pieces \emph{only} at an \code{await}.

That is the whole model, and two consequences fall out of it immediately.

\textbf{You never need a lock to protect a critical section that contains no
\code{await}.} Between two await points your code is atomic with respect to
every other coroutine, because none of them can run. This is a genuine
simplification over threaded code and it is worth internalising early.

\textbf{Code that never awaits never yields.} A coroutine that computes for two
seconds without awaiting holds the loop for two seconds, and every other request
in the process waits. There is no pre-emption. The scheduler is cooperative,
which is a polite way of saying it cannot help you.

\mermaidfig{async-event-loop}{%
  Two coroutines interleaving. Neither runs faster than it would alone; they
  overlap their \emph{waiting}, which is the only thing \code{asyncio}
  optimises.}{%
  Sequence diagram of two coroutines interleaving at await points.}

\begin{dotnetbox}
Closest analogue: a single-threaded \code{SynchronizationContext} that never
posts to a thread pool. The instinct that \code{await} might resume you on a
different thread does not apply --- in a normal \code{asyncio} program there is
only ever one thread running your coroutines, so the whole category of ``which
thread am I on'' questions disappears, and with it \code{ConfigureAwait}.
\end{dotnetbox}

% ============================================================
\section{The GIL, honestly}
\label{sec:gil}
\index{GIL}
\index{concurrency!versus parallelism}

The Global Interpreter Lock permits one thread at a time to execute Python
bytecode. It is the most discussed and least usefully understood thing about the
language, so here it is as a measurement rather than an argument.

The workload below is pure arithmetic --- no I/O, nothing that releases the lock.
One unit of it takes 0.29\,s on its own. Four units are then run three ways, on a
four-core machine:

\begin{center}
\small
\begin{tabularx}{\linewidth}{@{}lrX@{}}
\toprule
\textbf{Four CPU-bound units} & \textbf{Wall clock} & \textbf{Relative to one unit} \\
\midrule
One unit, alone & 0.29\,s & 1.0$\times$ --- the baseline \\
Four threads & 1.09\,s & 3.7$\times$ --- essentially serialised \\
Four processes & 0.29\,s & 1.0$\times$ --- genuinely parallel \\
\bottomrule
\end{tabularx}
\end{center}

\noindent
Four threads took 3.7 times as long as one unit. Perfect serialisation would be
4.0. Threads bought approximately nothing, because only one of them could
execute bytecode at a time. Four processes took exactly as long as a single unit,
because each process holds its own interpreter and its own lock.

\begin{note}
The rule that follows is short. \textbf{I/O-bound work} --- model calls, HTTP,
databases, file reads --- is what \code{asyncio} is for, and it scales
beautifully, as \S\ref{sec:measurement} shows. \textbf{CPU-bound work} --- local
embeddings, parsing a large PDF, compression --- needs \emph{processes}. Threads
are for blocking I/O in a library that has no async version, and for nothing
else.

Blocking I/O calls release the GIL while they wait, which is precisely why a
thread is the right tool there and the wrong tool for arithmetic.
\end{note}

\begin{versionbox}
Python 3.13 introduced an experimental free-threaded build without the GIL. It
is worth knowing about and it is not yet a production answer for this workload:
single-threaded performance is lower, and much of the ecosystem's C extension
surface has not been validated against it. Mention it as awareness of where the
language is going, not as something you have deployed.
\end{versionbox}

% ============================================================
\section{Coroutines, tasks and futures}
\label{sec:cold-coroutine}
\index{coroutine!cold}
\index{asyncio!create\_task}

This is the section that pays for the chapter.

\begin{python}[caption={A coroutine does not start when you create it}, label={lst:cold}]
async def work(tag, out):
    out.append(tag)          # a visible side effect, on the first line
    return tag


async def main():
    out = []
    coro = work("coro", out)
    print("after creating the coroutine :", out)
    await asyncio.sleep(0)               # let the loop run anything it can
    print("after a loop tick            :", out)
    await coro
    print("after awaiting               :", out)
\end{python}

\begin{console}
after creating the coroutine : []
after a loop tick            : []
after awaiting               : ['coro']
\end{console}

Calling \code{work(...)} produced an object and ran none of the function. Not one
line. The side effect on the first line of the body did not happen, and a full
turn of the event loop did not make it happen. Only \code{await} did.

\begin{dotnetbox}
This is the largest behavioural difference from .NET, and it runs the opposite
way to your instincts. \code{async Task<T> Foo()} in C\# begins executing
synchronously the moment you call it, up to its first incomplete await; the
\code{Task} you get back is usually already in flight. In Python nothing is in
flight until you say so.

The practical consequence: in C\#, forgetting to await a \code{Task} still
performs the work and merely loses the result. In Python, forgetting to await a
coroutine performs \emph{no work at all}, and warns you at garbage-collection
time if you happen to be looking.
\end{dotnetbox}

\subsection{What \texorpdfstring{\code{create\_task}}{create-task} actually does}
\label{sec:create-task}

The usual advice is that \code{create\_task} ``starts it immediately, in the
background''. That is close enough to be useful and wrong enough to mislead when
you are reasoning about ordering:

\begin{python}[caption={\code{create\_task} schedules; it does not execute}, label={lst:create-task}]
out = []
t = asyncio.create_task(work("task", out))
print("after create_task, before any tick :", out)
await asyncio.sleep(0)
print("after one loop tick                :", out)
\end{python}

\begin{console}
after create_task, before any tick : []
after one loop tick                : ['task']
\end{console}

It hands the coroutine to the loop and returns immediately. The body still does
not run until the current coroutine yields --- which here happens at the
\code{sleep(0)}. The difference from a bare coroutine is not that it runs
sooner. It is that it will run \emph{without anyone awaiting it}.

\mermaidfig{async-cold-coroutine}{%
  The three states worth distinguishing, and the consequence at the bottom that
  costs people the most time.}{%
  Coroutine versus task lifecycle, and the sequential-loop trap.}

\subsection{The consequence}
\label{sec:sequential-trap}

\begin{python}[caption={A queue wearing the costume of concurrency}, label={lst:seq-trap}]
# Sequential. Each iteration waits for the previous one to finish.
for url in urls:
    await fetch(url)

# Concurrent. All of them are in flight before any of them completes.
await asyncio.gather(*(fetch(url) for url in urls))
\end{python}

The first form is not a subtly worse version of the second. It is twenty times
slower for twenty URLs, and \S\ref{sec:measurement} measures exactly that.

\begin{warning}
This is the most common performance defect in async Python, and it is invisible
in review because both forms read as ``asynchronous''. The tell is \code{await}
\emph{inside} a loop body over independent work. Awaiting in a loop is correct
when each iteration depends on the last; it is a defect when they do not.
\end{warning}

% ============================================================
\section{Four ways to run things together}
\label{sec:four-ways}
\index{asyncio!gather}
\index{asyncio!TaskGroup}
\index{structured concurrency}

Four primitives with four different failure semantics. The differences only show
up when something raises, which is exactly when you need to already know them.

\begin{center}
\small
\begin{tabularx}{\linewidth}{@{}lX@{}}
\toprule
\textbf{Primitive} & \textbf{On the first exception} \\
\midrule
\code{gather} & Propagates it. Siblings are \emph{not} cancelled and keep
running. \\[4pt]
\code{gather} with \code{return\_exceptions} & Nothing propagates. Exceptions
arrive as ordinary items in the result list. \\[4pt]
\code{TaskGroup} & Cancels every sibling, waits for all of them, then raises an
\code{ExceptionGroup}. \\[4pt]
\code{as\_completed} & Yields results in completion order; each exception
surfaces when you await that result. \\[4pt]
\code{wait} & Returns \code{(done, pending)} and raises nothing. You inspect and
clean up yourself. \\
\bottomrule
\end{tabularx}
\end{center}

\subsection{The difference that matters}
\label{sec:gather-vs-taskgroup}

The \code{gather} row says siblings keep running. That is worth seeing, because
``keeps running'' means ``after your error handler has already run'':

\begin{python}[caption={Same failure, two primitives, two outcomes}, label={lst:gather-vs-tg}]
async def boom(log):
    await asyncio.sleep(0.05)
    log.append("boom raised")
    raise ValueError("boom")


async def slow(name, delay, log):
    try:
        await asyncio.sleep(delay)
        log.append(f"{name} finished")
    except asyncio.CancelledError:
        log.append(f"{name} CANCELLED")
        raise


# --- with gather ---
try:
    await asyncio.gather(boom(log), slow("sibling", 0.3, log))
except ValueError as e:
    log.append(f"gather propagated {e!r}")
await asyncio.sleep(0.4)        # wait around and watch what happens

# --- with TaskGroup ---
try:
    async with asyncio.TaskGroup() as tg:
        tg.create_task(boom(log))
        tg.create_task(slow("sibling", 0.3, log))
except* ValueError as eg:
    log.append(f"ExceptionGroup with {len(eg.exceptions)}")
\end{python}

\begin{console}
=== gather ===
   boom raised
   gather propagated ValueError('boom')
   sibling finished          <-- after the handler. Nobody is listening.
=== TaskGroup ===
   boom raised
   sibling CANCELLED
   ExceptionGroup with 1
\end{console}

Under \code{gather}, the sibling completed \emph{after} the exception had been
caught and handled. In a request handler that is a task still holding a database
connection and still writing to state belonging to a request that has already
failed and returned. Under \code{TaskGroup} it was cancelled, and the block did
not exit until it had finished being cancelled.

\mermaidfig{async-gather-vs-taskgroup}{%
  The orphaning, drawn. The branch marked in red is the one that causes
  production incidents, because it no longer has any relationship to the caller.}{%
  What happens to sibling tasks when one raises, under each primitive.}

\subsection{Structured concurrency}
\label{sec:structured-concurrency}

\code{TaskGroup} is not a nicer \code{gather}. It is an instance of a design
principle worth adopting deliberately: \textbf{a task may not outlive the block
that created it.}

That single rule eliminates a whole category of bug. No orphans. No leaked work
after a failure. No task quietly holding a connection open after its request has
gone. When control leaves the \code{async with}, every task started inside it has
finished --- successfully, or by being cancelled and having its cleanup run.

\begin{note}
Prefer \code{TaskGroup} in new code. Reach for \code{gather} in two cases: when
partial failure is acceptable and you want \code{return\_exceptions} semantics so
you can inspect each result; and when you need results as an ordered list
corresponding to the inputs, which \code{gather} gives you and \code{TaskGroup}
does not.

\code{TaskGroup} and \code{except*} require Python 3.11. This book assumes it.
\end{note}

\subsection{Completion order, and the low-level option}
\label{sec:as-completed}

\code{as\_completed} yields results as they finish, not in submission order:

\begin{python}[caption={Completion order, not submission order}, label={lst:as-completed}]
for fut in asyncio.as_completed([took("slow", 0.3), took("fast", 0.05),
                                 took("mid", 0.15)]):
    print(await fut)
\end{python}

\begin{console}
fast
mid
slow
\end{console}

That is the one to reach for when you want to start showing results before the
slowest branch finishes --- a retrieval fan-out where the first useful document
should appear immediately rather than after the slow shard replies.

\code{wait} is the low-level primitive. It raises nothing, returns two sets, and
leaves cancellation and cleanup to you. It is right when you genuinely want
``give me whatever is done after two seconds and let me decide about the rest'',
and it is the wrong default because that decision is then one you must remember
to make.

% ============================================================
\section{The measurement}
\label{sec:measurement}
\index{measurement!concurrency}

The claim under test: an I/O-bound workload interleaves almost perfectly under
\code{asyncio}, and the ceiling is whatever cap you impose rather than anything
about the machine.

\textbf{Method.} A real HTTP server on \code{localhost}, hand-rolled directly on
the standard library's stream-server API so that what is measured is
unambiguous: read a request, wait 100\,ms, reply. Twenty requests per trial over
one shared \code{httpx.AsyncClient} with a warmed connection pool. One warm-up
trial discarded, then ten timed trials; the median is reported. The script is
\code{code/ch03/bench\_concurrency.py} and runs in about thirty seconds.

Deliberately \emph{not} measured with \code{asyncio.sleep} on the client side.
Sleeping demonstrates that the scheduler can schedule; it says nothing about
sockets, and it would flatter the result.

\begin{center}
\small
\begin{tabularx}{\linewidth}{@{}lrrrX@{}}
\toprule
\textbf{Strategy} & \textbf{Median} & \textbf{Min} & \textbf{Max} & \textbf{Round trips} \\
\midrule
Sequential & 2.039\,s & 2.037 & 2.041 & 20.4$\times$ \\
\code{gather} & 0.126\,s & 0.124 & 0.132 & 1.3$\times$ \\
\code{TaskGroup} & 0.125\,s & 0.124 & 0.128 & 1.2$\times$ \\
Capped at 5 & 0.421\,s & 0.420 & 0.424 & 4.2$\times$ \\
\bottomrule
\end{tabularx}
\end{center}

\noindent
Twenty requests, 100\,ms server delay. ``Round trips'' is the median divided by
one round trip, which is the useful unit because it is machine-independent.

Four things are worth drawing out.

\textbf{Sequential cost 20.4 round trips for 20 requests.} It is a queue. The
theoretical floor is 20 and it hit 20.4, so essentially all of the elapsed time
was server delay and almost none of it was overhead.

\textbf{Concurrent cost 1.3 round trips.} The floor is 1. Twenty simultaneous
requests finished in very slightly more than the time for one, which is what
``I/O-bound'' means: the process was not working, it was waiting, and waits
overlap for free. That is a \textbf{16$\times$ improvement} from a one-line
change.

\textbf{\code{TaskGroup} and \code{gather} are indistinguishable on speed} ---
0.125 against 0.126\,s. The choice between them is a failure-semantics decision
(\S\ref{sec:gather-vs-taskgroup}) and never a performance one. Anyone choosing
\code{gather} for speed is choosing on a difference that does not exist.

\textbf{The cap sets the ceiling, exactly.} Capped at five, twenty requests took
4.2 round trips against a theoretical floor of four. The concurrency limit does
precisely what it says, and its cost is arithmetic you can predict before you
deploy. That matters because \S\ref{sec:semaphore} argues you should almost
always impose one.

\begin{note}
The spread is tiny --- 0.124 to 0.132\,s across ten trials --- because
\code{localhost} has no network variance. Against a real provider the spread is
the interesting part rather than the median, and Chapter~\ref{ch:serving}
measures under conditions where that is true. Take the ratios from this
experiment, which are structural, rather than the absolute numbers, which are
not.
\end{note}

% ============================================================
\section{Limiting concurrency}
\label{sec:semaphore}
\index{asyncio!Semaphore}
\index{rate limiting}

\S\ref{sec:measurement} makes unbounded fan-out look free. Against a model
provider it is not.

Twenty concurrent calls to a provider is twenty concurrent calls against a rate
limit you do not control. Two thousand is a wave of \code{429}s, a retry storm,
and --- because each in-flight request holds its request and response in memory
--- a plausible route to an out-of-memory kill. The primitive is one line:

\begin{python}[caption={Capping concurrency}, label={lst:semaphore}]
sem = asyncio.Semaphore(5)


async def guarded(prompt):
    async with sem:
        return await call_model(prompt)


async with asyncio.TaskGroup() as tg:
    for p in prompts:
        tg.create_task(guarded(p))
\end{python}

\begin{dotnetbox}
\code{asyncio.Semaphore} is \code{SemaphoreSlim}, with the same role and
essentially the same shape --- \code{async with} in place of
\code{await WaitAsync()} plus a \code{finally} that releases. The context-manager
form means you cannot forget the release, which is the bug this construct
attracts in both languages.
\end{dotnetbox}

\begin{warning}
A semaphore is per-process. The moment you run more than one worker --- and
Chapter~\ref{ch:serving} runs several --- your effective concurrency is the
limit multiplied by the number of workers, which is neither the number you
configured nor the number the provider's rate limit is expressed in.
Chapter~\ref{ch:serving} is where that gets resolved. For now: know the number is
per-process, and size it accordingly.
\end{warning}

% ============================================================
\section{Queues and backpressure}
\label{sec:queues}
\index{asyncio!Queue}
\index{backpressure}

\code{asyncio.Queue} is the producer/consumer primitive, and its most important
argument is the one people leave off.

\begin{python}[caption={A bounded queue is backpressure}, label={lst:queue}]
q = asyncio.Queue(maxsize=2)
await q.put(1)
await q.put(2)
await q.put(3)     # blocks here until a consumer takes something
\end{python}

With \code{maxsize} set, \code{put} blocks when the queue is full. That is not an
inconvenience --- it is the mechanism by which a slow consumer slows its producer
down, which is the only thing standing between you and unbounded memory growth.

Leave \code{maxsize} off and the queue is unbounded. A producer faster than its
consumer then converts a throughput mismatch into a memory leak with a delay
fuse. Chapter~\ref{ch:serving} meets the production instance: a browser on hotel
wifi consuming a token stream more slowly than the model produces it.

\begin{dotnetbox}
\code{asyncio.Queue(maxsize=n)} is \code{Channel.CreateBounded<T>(n)}, and
leaving it off is \code{Channel.CreateUnbounded<T>()}. Same trade, same reasons,
same rule: bounded unless you can prove the producer cannot outrun the consumer.
\end{dotnetbox}

The other synchronisation primitives --- \code{Lock}, \code{Event},
\code{Condition}, \code{Barrier} --- exist and mirror their threading
counterparts. \code{Lock} is needed far less often than instinct suggests,
because of \S\ref{sec:event-loop}: code between awaits is already atomic. Reach
for it when a critical section \emph{does} contain an await, which is the only
case where another coroutine can interleave.

% ============================================================
\section{Async generators and context managers}
\label{sec:async-generators}
\index{asyncio!async generator}
\index{asyncio!aclosing}

Streaming is an async generator, and every streaming interface in this book is
one:

\begin{python}[caption={The shape of every streaming API here}, label={lst:agen}]
async def stream_tokens(prompt):
    async for chunk in model.astream(prompt):
        yield chunk.text
\end{python}

\begin{dotnetbox}
\code{async def ... yield} is \code{IAsyncEnumerable<T>} with
\code{yield return}; \code{async for} is \code{await foreach}. \code{async with}
is \code{await using}, and the protocol is a pair of dunder methods rather than
\code{IAsyncDisposable}.
\end{dotnetbox}

The trap is abandonment. Break out of an \code{async for} early and the generator
is left suspended at its \code{yield}, holding whatever it had open --- an HTTP
response, a cursor, a connection --- until the garbage collector reaches it,
which is not a schedule you control:

\begin{python}[caption={Closing a generator you abandon}, label={lst:aclosing}]
from contextlib import aclosing

async with aclosing(stream_tokens(prompt)) as stream:
    async for token in stream:
        if user_disconnected():
            break          # aclosing runs the generator's cleanup on exit
\end{python}

This is not hypothetical for a streaming service: a client disconnecting
mid-response is the normal case rather than the exceptional one, and
Chapter~\ref{ch:cancellation} takes it further.

% ============================================================
\section{Blocking code}
\label{sec:blocking}
\index{asyncio!to\_thread}
\index{blocking call}

The failure mode this chapter has been building towards, measured.

Ten healthy coroutines each await 50\,ms of I/O. Then one more coroutine is
added, which calls a blocking library function taking 500\,ms. The question is
what happens to the other ten.

\begin{center}
\small
\begin{tabularx}{\linewidth}{@{}lrX@{}}
\toprule
\textbf{Scenario} & \textbf{Worst healthy request} & \\
\midrule
No offender & 50.3\,ms & the baseline \\
Blocking offender & 500.4\,ms & \textbf{10$\times$ worse} \\
Offender wrapped in \code{to\_thread} & 50.9\,ms & back to baseline \\
\bottomrule
\end{tabularx}
\end{center}

\noindent
One coroutine, one blocking call, and every unrelated request in the process got
ten times slower. Nothing errored. Nothing logged. The requests were not
competing for a resource --- they were simply unable to run, because the loop was
inside \code{time.sleep} and there is no pre-emption.

\mermaidfig{async-blocking-call}{%
  The blocking call, with the measured numbers on it. Note what the diagram does
  not contain: an error, a warning, or any signal at all.}{%
  How one blocking call degrades every unrelated request.}

\begin{warning}
This is the most damaging single defect in async Python and the hardest to find
after the fact, because the symptom appears on requests that have nothing to do
with the cause. Latency rises across the board, traces show slow spans
everywhere, and nothing points at the offending line.

Chapter~\ref{ch:cancellation} covers finding one in a running process. Avoiding
it is cheaper. The usual offenders: \code{requests}, \code{time.sleep},
\code{psycopg2}, any synchronous database driver, file I/O on a large file, and
JSON parsing of a very large payload.
\end{warning}

\subsection{The two escape hatches}
\label{sec:to-thread}

\begin{python}[caption={Getting blocking work off the loop}, label={lst:to-thread}]
# Blocking I/O in a library with no async version: a thread.
text = await asyncio.to_thread(pdf_extract_text, path)

# CPU-bound work: a process. A thread would not help -- see the GIL table.
loop = asyncio.get_running_loop()
with ProcessPoolExecutor() as pool:
    vectors = await loop.run_in_executor(pool, embed_locally, chunks)
\end{python}

\begin{center}
\small
\begin{tabularx}{\linewidth}{@{}lX@{}}
\toprule
\textbf{Use} & \textbf{For} \\
\midrule
A native async library & Always, when one exists. \code{httpx} over
\code{requests}; \code{psycopg} 3 or \code{asyncpg} over \code{psycopg2}. \\[4pt]
\code{asyncio.to\_thread} & Blocking \emph{I/O} with no async alternative. The
call releases the GIL while it waits. \\[4pt]
\code{ProcessPoolExecutor} & CPU-bound work. Threads are measurably useless here
(\S\ref{sec:gil}). \\
\bottomrule
\end{tabularx}
\end{center}

\begin{warning}
\code{to\_thread} is not free and the pool is not infinite. It runs on the
default executor, whose worker count is bounded; saturate it and calls queue,
which reproduces the original problem one level down and with a less obvious
symptom. It is a patch for a library you cannot replace, not an architecture.
\end{warning}

% ============================================================
\section{Why this chapter is in a book about agents}
\label{sec:why-async-agents}

Because an agent is the purest I/O-bound workload most engineers will ever
deploy.

A single turn is: send a prompt, wait one to thirty seconds, receive tool calls,
wait on an API or a database, send the results back, wait again. Better than
ninety per cent of wall-clock time is spent waiting on somebody else's network.
There is almost no computation in an agent --- the computation happens on the
provider's hardware, and you are billed for it separately.

\S\ref{sec:measurement} put a number on what that means: a 16$\times$ difference
between the obvious implementation and the correct one, on a workload whose shape
is exactly an agent's. Without \code{asyncio}:

\begin{itemize}[leftmargin=1.4em,nosep]
\item one user occupies one worker for the whole of a twenty-second turn;
\item five tools that could be called at once are called one after another;
\item token streaming is not really available;
\item infrastructure cost rises linearly with concurrent users, for a process
      that is idle almost all the time.
\end{itemize}

\begin{note}
Worth having ready as a sentence: \emph{an agent workload is dominated by
provider latency, so one process can hold hundreds of concurrent conversations
while spending almost all of its time idle --- and the single thing that destroys
that property is one blocking call, which costs an order of magnitude and
produces no error.} Both halves of that are measured in this chapter rather than
assumed.
\end{note}

% ============================================================
\section{What to take away}
\label{sec:ch3-summary}

\begin{itemize}[leftmargin=1.4em]
\item A coroutine is cold. It runs when awaited, not when created.
      \code{create\_task} schedules it; the body still waits for the next yield.
\item \code{await} inside a loop over independent work is a queue. Measured:
      20.4 round trips against 1.3.
\item \code{gather} leaves siblings running after an exception has been handled.
      \code{TaskGroup} cancels them and waits. They are identical on speed, so
      the choice is always about failure semantics.
\item Structured concurrency: a task may not outlive the block that created it.
\item Cap fan-out with a \code{Semaphore}. The cap sets the ceiling exactly ---
      measured at 4.2 round trips for a cap of five over twenty requests.
\item Bound your queues. An unbounded queue with a slow consumer is a memory leak
      with a delay fuse.
\item Close async generators you abandon, with \code{aclosing}.
\item One blocking call degrades every unrelated request by 10$\times$ and emits
      no signal. Native async library first, \code{to\_thread} for blocking I/O,
      a process for CPU work.
\end{itemize}

\begin{exercisebox}
\textbf{1.} Run \code{code/ch03/bench\_concurrency.py}. Then set the server delay
to 1\,ms and run it again. The ratios collapse --- work out why before reading
on, because the answer is the definition of ``I/O-bound''.

\textbf{2.} Reproduce the orphaned task. Use \code{gather} with one failing and
one slow coroutine, catch the exception, and print from the slow one when it
finally finishes. Then convert it to a \code{TaskGroup} and watch the print
disappear.

\textbf{3.} Take the blocking-call measurement in \S\ref{sec:blocking} and add a
second, third and fourth blocking coroutine. Latency for the healthy requests
should rise in steps you can predict exactly. Predict them first.

\textbf{4.} Write a producer/consumer with \code{Queue(maxsize=1)} where the
producer is ten times faster than the consumer. Log queue depth every 100\,ms.
Then remove \code{maxsize} and watch the depth --- and the process's memory ---
instead.

\textbf{5.} Find a blocking call in code you already own. Every non-trivial
Python service has one. \code{requests}, \code{time.sleep} and a synchronous
database driver are where to look first.
\end{exercisebox}

\begin{projectbox}
\textbf{Stage 01 --- the concurrency playground.} No agent yet: the primitives
every later stage depends on, written once, where they can be understood.

A capped fan-out under \code{TaskGroup}; a bounded producer/consumer queue; and
the benchmark from \S\ref{sec:measurement} as a reproducible script.
Chapter~\ref{ch:cancellation} completes the stage with graceful shutdown.

The test contract is in \code{stages/01-async/README.md}. Its most important
entry: concurrency must never exceed the cap, asserted on \emph{observed peak}
rather than on wall-clock time. A timing-based assertion for this passes on a
fast machine and fails in CI, which teaches the reader the wrong lesson about
their own code.
\end{projectbox}
